\documentclass[conference]{IEEEtran}
\IEEEoverridecommandlockouts

\usepackage{cite}
\usepackage{amsmath,amssymb,amsfonts}
\usepackage{algorithmic}
\usepackage{graphicx}
\usepackage{textcomp}
\usepackage{xcolor}
\usepackage{booktabs}
\usepackage[hidelinks]{hyperref}

\def\BibTeX{{\rm B\kern-.05em{\sc i\kern-.025em b}\kern-.08em
    T\kern-.1667em\lower.7ex\hbox{E}\kern-.125emX}}

\begin{document}

\title{Paper Title: Subtitle Goes Here}

\author{
  \IEEEauthorblockN{1\textsuperscript{st} Given Name Surname}
  \IEEEauthorblockA{
    \textit{Department of Computer Science} \\
    \textit{University / Organization} \\
    City, Country \\
    email@example.com
  }
  \and
  \IEEEauthorblockN{2\textsuperscript{nd} Given Name Surname}
  \IEEEauthorblockA{
    \textit{Department of Computer Science} \\
    \textit{University / Organization} \\
    City, Country \\
    email@example.com
  }
}

\maketitle

\begin{abstract}
This document is a minimal IEEE conference paper template based on \texttt{IEEEtran}. Replace this text with a concise summary of the problem, methodology, key results, and conclusions (typically 150--250 words).
\end{abstract}

\begin{IEEEkeywords}
component, formatting, style, styling, insert
\end{IEEEkeywords}

\section{Introduction}
Introduce the motivation, problem statement, and contributions of the work. Cite prior work using~\cite{example_conf} or~\cite{example_journal}.

\section{Background and Related Work}
Summarize the technical background and compare with closely related studies.

\section{Design and Implementation}
Describe the architecture, algorithms, and implementation details.

\begin{table}[t]
  \centering
  \caption{Example Summary Table}
  \label{tab:example}
  \begin{tabular}{@{}lcc@{}}
    \toprule
    Method & Metric A & Metric B \\
    \midrule
    Baseline & 1.00 & 10.5 \\
    Proposed & \textbf{0.42} & \textbf{24.8} \\
    \bottomrule
  \end{tabular}
\end{table}

\section{Evaluation}
Present the experimental setup, workloads, and quantitative results.

% Example figure block:
% \begin{figure}[t]
%   \centering
%   \includegraphics[width=\linewidth]{figures/example.pdf}
%   \caption{Example figure caption.}
%   \label{fig:example}
% \end{figure}

\section{Conclusion}
Summarize the findings, limitations, and future work.

\section*{Acknowledgment}
Optional acknowledgment section.

\bibliographystyle{IEEEtran}
\bibliography{references}

\end{document}
